% !TeX root = ../../Book5.tex
% 纲要标题：### 13.1 不变双线性型
% 纲要位置：工作记录/计划纲要.md，第 4143 行。

\section{不变双线性型}
\label{b5:sec:1301}

% 纲要范围：
%
% * Killing 型
% * 不变性及根空间正交的动机
% * 简单李代数上的形式
%

上一章用上三角矩阵识别可解性。本章通过伴随算子的乘积取迹，构造一个双线性型。Cartan 判据将证明，在有限维复李代数中，这个型非退化恰好等价于李代数半单。

\subsection{Killing 型}
\label{b5:subsec:1301:01}

\begin{definition}\label{b5:def:killing-form}
设 \(\mathfrak g\) 是域 \(k\) 上的有限维李代数。映射
\[
\kappa_{\mathfrak g}(x,y)
=\operatorname{tr}_{\mathfrak g}
\Bigl(\operatorname{ad}x\circ\operatorname{ad}y\Bigr),
\qquad x,y\in\mathfrak g,
\]
称为 \(\mathfrak g\) 的\term[Killing-xing]{Killing 型}[Killing form]。更一般地，若 \(V\) 是有限维 \(k\) 向量空间，表示
\(\rho:\mathfrak g\rightarrow\mathfrak{gl}(V)\) 给出\term[ji-xing]{迹型}[trace form]
\(\beta_\rho(x,y)=\operatorname{tr}_V\Bigl(\rho(x)\rho(y)\Bigr)\)。
\symidx{\(\kappa_{\mathfrak g}\)}
\end{definition}

\begin{example}\label{b5:exm:killing-sln}
对 \(X,Y\in\mathfrak{sl}_n(\mathbb C)\)，
\begin{equation}\label{b5:eq:killing-sln}
\kappa_{\mathfrak{sl}_n}(X,Y)=2n\operatorname{tr}(XY).
\end{equation}
为核验系数，在 \(M_n(\mathbb C)\) 上记 \(L_X(T)=XT\)、\(R_Y(T)=TY\)。矩阵单位基给出
\[
\operatorname{tr}(L_A)=\operatorname{tr}(R_A)=n\operatorname{tr}A,
\qquad
\operatorname{tr}(L_XR_Y)=\operatorname{tr}X\,\operatorname{tr}Y.
\]
最后一式中 \(E_{ij}\) 的自身系数为 \(X_{ii}Y_{jj}\)，求和即可。展开
\((L_X-R_X)(L_Y-R_Y)\)，得
\[
\kappa_{\mathfrak{gl}_n}(X,Y)
=2n\operatorname{tr}(XY)-2\operatorname{tr}X\,\operatorname{tr}Y.
\]
又因 \(\mathfrak{gl}_n=\mathfrak{sl}_n\oplus\mathbb CI\)，伴随作用在标量直线上为零，便得到\eqref{b5:eq:killing-sln}。特别在 \(n=2\) 时，对标准的 \(e=E_{12}\)、\(f=E_{21}\)、\(h=E_{11}-E_{22}\)，有
\[
\kappa(h,h)=8,\qquad \kappa(e,f)=4,\qquad
\kappa(e,e)=\kappa(f,f)=\kappa(h,e)=\kappa(h,f)=0.
\]
\end{example}

\subsection{不变性与根空间正交}
\label{b5:subsec:1301:02}

\begin{definition}\label{b5:def:invariant-bilinear-lie}
设 \(\mathfrak g\) 是域 \(k\) 上的李代数，\(B:\mathfrak g\times\mathfrak g\rightarrow k\) 是双线性型。若
\[
B\Bigl([x,y],z\Bigr)=B\Bigl(x,[y,z]\Bigr)\qquad(x,y,z\in\mathfrak g),
\]
则 \(B\) 称为一个\term[bubian-shuangxianxing-xing]{不变双线性型}[invariant bilinear form]。
\end{definition}

\begin{proposition}\label{b5:prop:invariant-trace-form}
设 \(\mathfrak g\) 为域 \(k\) 上的有限维李代数。它的每个有限维表示的迹型均为对称不变双线性型，因此其 Killing 型对称且不变。对任意对称不变双线性型 \(B:\mathfrak g\times\mathfrak g\to k\) 和理想 \(I\subseteq\mathfrak g\)，正交补
\[
I^\perp=\Bigl\{\,x\in\mathfrak g\mid B(x,I)=0\,\Bigr\}
\]
也是理想；特别地，型的根 \(\mathfrak g^\perp\) 是理想。
\end{proposition}
\begin{proof}
矩阵迹满足 \(\operatorname{tr}(AB)=\operatorname{tr}(BA)\)，并且
\[
\operatorname{tr}\Bigl([A,B]C\Bigr)=\operatorname{tr}\Bigl(A[B,C]\Bigr).
\]
对表示矩阵代入即得对称性与不变性。若 \(u\in I^\perp\)、\(x\in\mathfrak g\)、\(y\in I\)，则
\[
B\Bigl([x,u],y\Bigr)=-B\Bigl(u,[x,y]\Bigr)=0
\]
因为 \([x,y]\in I\)。于是 \([x,u]\in I^\perp\)。
\end{proof}

\begin{proposition}\label{b5:prop:weight-orthogonality-invariant}
设 \(\mathfrak g\) 是域 \(k\) 上的李代数，\(B:\mathfrak g\times\mathfrak g\to k\) 是对称不变双线性型，\(\mathfrak h\subseteq\mathfrak g\) 是子代数，\(\alpha,\beta\in\operatorname{Hom}_k(\mathfrak h,k)\)，\(x,y\in\mathfrak g\)。若
\[
[h,x]=\alpha(h)x,\qquad[h,y]=\beta(h)y
\quad(h\in\mathfrak h)
\]
且 \(\alpha+\beta\ne0\)，则 \(B(x,y)=0\)。
\end{proposition}
\begin{proof}
不变性给出
\[
0=B\Bigl([h,x],y\Bigr)+B\Bigl(x,[h,y]\Bigr)
=\Bigl(\alpha(h)+\beta(h)\Bigr)B(x,y).
\]
取使 \(\alpha(h)+\beta(h)\ne0\) 的 \(h\)，即可除去该非零标量。
\end{proof}

后面的根空间正交关系将直接使用这条命题。它目前只是一条关于共同特征向量的计算，不要求已经存在 Cartan 子代数或根系。

\begin{lemma}\label{b5:lem:killing-restriction-ideal}
设 \(\mathfrak g\) 是域 \(k\) 上的有限维李代数，\(I\) 是理想。对 \(x,y\in I\)，有
\[
\kappa_{\mathfrak g}(x,y)=\kappa_I(x,y).
\]
\end{lemma}
\begin{proof}
伴随算子 \(\operatorname{ad}x\) 把 \(\mathfrak g\) 送入 \(I\)，在 \(\mathfrak g/I\) 上诱导零作用。将 \(I\) 的基扩充为 \(\mathfrak g\) 的基，算子乘积的矩阵呈分块上三角形，商块为零，故其全空间迹等于限制到 \(I\) 的迹。
\end{proof}

\subsection{简单李代数上的形式}
\label{b5:subsec:1301:03}

\begin{proposition}\label{b5:prop:simple-invariant-form-unique}
设 \(\mathfrak g\) 是非零有限维复简单李代数。则每个非零对称不变双线性型均非退化；这种型所成空间的维数至多为一。
\end{proposition}
\begin{proof}
由\cref{b5:prop:invariant-trace-form}，一个不变型的根是理想。非零型的根不可能是整个 \(\mathfrak g\)，简单性便迫使其根为零。

固定一个非零型 \(B_0\)。对另一个型 \(B\)，非退化性唯一确定线性算子 \(T\)，满足 \(B(x,y)=B_0(Tx,y)\)。两个型的不变性说明
\[
B_0\Bigl(T[a,x],y\Bigr)=B_0\Bigl([a,Tx],y\Bigr)
\]
对所有 \(y\) 成立，因而 \(T\) 与全部 \(\operatorname{ad}a\) 交换。复数域上 \(T\) 有特征值 \(c\)，非零空间 \(\ker(T-cI)\) 是伴随子模，即 \(\mathfrak g\) 的理想。因此它等于 \(\mathfrak g\)，\(T=cI\)，于是 \(B=cB_0\)。
\end{proof}

这时还不能仅凭此命题宣称 Killing 型非零：它只说明“若已有一个非零不变型，则唯一到标量”。下一节的判据将保证复简单李代数确实具有非退化 Killing 型，从而把“至多一维”加强为“恰好一维”。
